\begin{pastexam}[label={px:ust-112-ee-ds-12}]{臺聯大電機類 112 學年度 第 12 題}
（複選；3 分，每一選項單獨計分，不倒扣）

Use a stack structure with only push and pop operations. An array with 4 items \{A,B,C,D\} is
processed by push(A), push(B), push(C), push(D) sequentially with all possible legal pop() to remove
the item on top of the stack at any time. For example, we can execute push(A), push(B), push(C),
push(D), and then pop(), pop(), pop(), pop() to reverse the order to get D, C, B, A pop
sequentially. What of the following are not possible pop sequences using a stack?

(A)~ACBD \qquad (B)~CDBA \qquad (C)~CDAB \qquad (D)~BDCA \qquad (E)~BDAC
\end{pastexam}

\begin{pastsoln}
答案：(C)、(E)。

模擬：照輸出序列一個一個要，要的元素如果不在堆疊頂端，就把還沒 push 的依序 push 進去，直到
它出現在頂端再 pop；如果它已經在堆疊裡、但不在頂端，就不可能（資結書第~6.1.5~節）。

\begin{itemize}
\item (A) ACBD：push A、pop A；push B、push C、pop C；pop B；push D、pop D。 可能。
\item (B) CDBA：push A、B、C，pop C；push D、pop D；堆疊剩 A、B（B 在頂），pop B、pop A。 可能。
\item (C) CDAB：push A、B、C，pop C；push D、pop D；這時堆疊是 A、B，頂端是 B，要的是 A，
A 被 B 壓住。 \textbf{不可能}。
\item (D) BDCA：push A、B，pop B；push C、D，pop D；pop C；pop A。 可能。
\item (E) BDAC：push A、B，pop B；push C、D，pop D；堆疊是 A、C，頂端是 C，要的是 A。
\textbf{不可能}。
\end{itemize}

也可以用資結書定理~6.1 判斷：依序 push 的輸出序列不可能，若且唯若裡面有三個數
（照 push 順序叫小、中、大）以「大、小、中」的順序出現。 (C) 的 D、A、B 和 (E) 的 D、A、C 都是
這種樣子；(A)(B)(D) 找不到。
\end{pastsoln}
